\documentclass[10pt,a4paper]{amsart}
\usepackage[margin=19mm]{geometry}
\usepackage{amsmath,amssymb,mathtools}
\usepackage[scr=rsfs,cal=euler]{mathalfa}
\usepackage{xfrac}
\usepackage[unicode,hidelinks,bookmarks=false]{hyperref}
\newcommand{\ie}{i.e.\ }
\newcommand{\cf}{cf.\ }
\newcommand{\resp}{respectively\ }
\newcommand{\Subsection}{\subsection}
\providecommand{\nobreakdash}{\nobreak\hbox{-}}
\setlength{\emergencystretch}{2em}
\multlinegap0pt
\usepackage{bm}
\usepackage{bbm}
\usepackage{dsfont}
\usepackage{xspace}
\usepackage{cancel}
\usepackage{mathtools}
\usepackage{amsthm}
\makeatletter
\@ifundefined{theorem}{\newtheorem{theorem}{Theorem}}{}
\@ifundefined{proposition}{\newtheorem{proposition}[theorem]{Proposition}}{}
\makeatother
\newcommand{\A}{\mathcal{A}}
\newcommand{\C}{\mathcal{C}}
\newcommand{\E}{\mathbb{E}}
\newcommand{\Q}{\mathcal{Q}}
\newcommand{\X}{\mathcal{X}}
\newcommand{\Y}{\mathcal{Y}}
\newcommand{\Z}{\mathbb{Z}}
\newcommand{\eq}[1]{\begin{align*}#1\end{align*}}
\newcommand{\e}{\epsilon}
\newcommand{\norm}[1]{\left\lVert#1\right\rVert}
\renewcommand{\P}{\mathcal{P}}
\renewcommand{\S}{\mathcal{S}}
\title[Upper and lower bounds on the subgeometric convergence of ad]{Upper and lower bounds on the subgeometric convergence of adaptive Markov chain Monte Carlo}
\author{Austin Brown, Jeffrey S. Rosenthal}
\date{Source: arXiv:2411.17084v3 (CC BY 4.0)}
\begin{document}
\maketitle

\noindent\emph{Source and licence.} Upper and lower bounds on the subgeometric convergence of adaptive Markov chain Monte Carlo. Version arXiv:2411.17084v3 (CC BY 4.0), distributed under the Creative Commons Attribution 4.0 International licence, as recorded in the arXiv licence metadata for this exact version and shown on the abstract page https://arxiv.org/abs/2411.17084v3. Full rights notice: licenses/arxiv-2411.17084-CC-BY-4.0.txt.\par\smallskip

\noindent\textit{Editorial scope.} Counted unit: the Proposition whose source label is \texttt{proposition:lb\_imh} in the article (source0.tex lines 853--864, source label \texttt{proposition:lb\_imh}), delivered together with the article's own complete original proof of it (source0.tex lines 866--932, one \texttt{proof} environment of 67 lines). No mathematics is written, repaired or re-derived: statement and proof are the article's own text line for line. Required context retained uncounted: theorem:lb\_with\_target\_tails (lines 198--221); eq:mhi (lines 843--848). Every definition, display and label of those lines is retained unchanged. Omitted from this package and not redistributed: 1--197, 222--842, 849--852, 865--865, 933--1496. Nothing inside a retained excerpt is omitted or altered: every retained body is delivered exactly as the article's lines are written, so no omission receipt of any kind accompanies this package. Citations: the retained excerpts carry no citation command at all, so no bibliography entry of the article is transcribed and no reference is redistributed. Reference adaptation: none was needed and none was made; every reference inside the delivered unit resolves inside it, so no unresolved reference or citation is produced. Printed slips kept literal: no printed word, symbol, hypothesis or number of the counted statement or of its proof was altered. Layout and class: the article's own class is \texttt{article} with packages including inputenc, geometry, natbib, amsthm, amsmath, amsfonts, amssymb, graphicx, setspace, authblk, hyperref; the wrapper is the lane's pinned \texttt{amsart} editorial wrapper at 10pt on A4, which restates the theorem-like environments the retained text needs and adds the article's own macro definitions that the retained text needs (\texttt{\textbackslash A}, \texttt{\textbackslash C}, \texttt{\textbackslash E}, \texttt{\textbackslash P}, \texttt{\textbackslash Q}, \texttt{\textbackslash S}, \texttt{\textbackslash X}, \texttt{\textbackslash Y}, \texttt{\textbackslash Z}, \texttt{\textbackslash e}, \texttt{\textbackslash eq}, \texttt{\textbackslash norm}), with extra wrapper packages (bm, bbm, dsfont, xspace, cancel, mathtools). The delivered numbering is the unit's own. Assets: no retained span contains a figure environment, a graphics-inclusion command, a PDF-page inclusion, a table or a TikZ picture; no image file is copied into this package, the package declares no asset dependency and no image file is redistributed with it. Licence: version arXiv:2411.17084v3 of this article is distributed under the Creative Commons Attribution 4.0 International licence, as recorded in the arXiv licence metadata for this exact version and shown on the abstract page https://arxiv.org/abs/2411.17084v3. A full rights notice is delivered at licenses/arxiv-2411.17084-CC-BY-4.0.txt. Characters: every retained excerpt is pure ASCII, so the delivered engine, XeLaTeX with its default Latin Modern text font, reports no missing character.
\begin{theorem}
\label{theorem:lb_with_target_tails}
Assume there is a Borel function $W : \X \to [1, \infty)$ and constants $C \in (0, 1]$ and $\kappa \in (0, \infty)$ where
\begin{align}
\pi( W \ge r ) \ge C r^{-\kappa}
\label{eq:target_tail_lb}
\end{align}
holds for all $r \ge 1$ and there is an $\alpha > \kappa$ and a concave function $\phi : (0, \infty) \to (0, \infty)$ such that
\begin{align}
&(\P_\gamma W^{\alpha} )(x) - W(x)^\alpha  \le \phi( W(x)^\alpha )
\label{eq:lb_drift}
\end{align}
holds for all $(x, \gamma) \in \X \times \Y$.
Then for all $t$,
\eq{
\inf_{\Q \in \S(\X, \Y) } \norm{ \A_\Q^{(t)}(\gamma_0, x_0, \cdot) - \pi }_{\text{TV}}
\ge \frac{M}{ \left( H_{W(x_0)^\alpha, \phi}^{-1}(t) \right)^{\frac{\kappa}{\alpha - \kappa}}}
}
where 
\begin{align}
M = C^{\frac{\alpha}{\alpha - \kappa}} \left[ ( \kappa/\alpha )^{\frac{\kappa}{\alpha - \kappa}} - ( \kappa/\alpha )^{\frac{\alpha}{\alpha - \kappa}} \right].
\label{eq:M}
\end{align}
\end{theorem}

\begin{align}
\P_\gamma(x, A)
= \int_A a_\gamma(x, y) \gamma \exp(-\gamma y) dy
+ \delta_{x}(A) R_\gamma(x)
\label{eq:mhi}
\end{align}

\begin{proposition}
\label{proposition:lb_imh}
Let $\A_{\Q}^{(t)}(\gamma_0, x_0, \cdot)$ be the marginal of the adaptive independent Metropolis-Hastings process at time $t \in \Z_+$ from \eqref{eq:mhi} with adaptation parameter set $(\gamma_*, \gamma^*)$ and initialized at $x_0, \gamma_0 \in (0, \infty) \times (\gamma_*, \gamma^*)$.
Then for any $\e \in (0, 1)$, there is a $\delta_{\e} \in (0, 1)$
\begin{align*}
\inf_{\Q \in \S([0, \infty), (\gamma_*, \gamma^*)) } \inf_{\xi \in \C\left[ \A_{\Q}^{(t)}(\gamma_0, x_0, \cdot), \pi \right] } \xi( \{ x, y \in \X \times \X : |y - x| > \delta_{\e} \})
\ge (1 - \e) \frac{M_*}{ \left( c_* t + \exp(\alpha x_0) \right)^{\frac{1}{\alpha - 1}} }
\end{align*}
for $\alpha \in (1, \gamma_*)$ with $M_* = \alpha^{-1/(\alpha - 1)} - \alpha^{-\alpha/(\alpha - 1)}$ and $c_{*} = \gamma^* / (\alpha - 1)
+ \gamma_* / (\gamma_* - \alpha)
+ \gamma^* - 1$.
\end{proposition}

\begin{proof}
Define $W(x) = \exp(x)$, and we have by a standard computation $\pi(W(x) \ge r) = r^{-1}$.
Assume $1 < \gamma_* < \gamma$.
For $\alpha < \gamma$, the identity holds
\eq{
&(\P_\gamma W^{\alpha})(x) - W^{\alpha}(x)
\\
&= \int_{[0, \infty)}  \exp(\alpha y) \left\{ 1 \wedge \left[ \exp\left[ (\gamma - 1) ( y - x ) \right] \right] \right\} \gamma \exp(-\gamma y) dy 
- \exp(\alpha x) \E\left( a_\gamma(x, Y) \right).
}
We also have for any $\alpha < \gamma$,
\begin{align}
&\int_{[0, \infty)}  \exp(\alpha y) \left\{ 1 \wedge \left[ \exp\left[ (\gamma - 1) ( y - x ) \right] \right] \right\} \gamma \exp(-\gamma y) dy
\nonumber
\\
&= 
\frac{\gamma \exp\left[ (1 - \gamma) x \right]}{\alpha - 1} \int_0^x (\alpha - 1) \exp\left[ (\alpha - 1) y \right] dy
\nonumber
\\
&\hspace{1.5em}+ \frac{\gamma}{\gamma - \alpha} \int_x^\infty (\gamma - \alpha) \exp\left[ - ( \gamma - \alpha) y \right] dy
\nonumber
\\
&= 
\left\{ \frac{\gamma}{\alpha - 1} 
+ \frac{\gamma}{\gamma - \alpha}
\right\} 
\exp\left[ -(\gamma - \alpha) x \right] 
- \frac{\gamma}{\alpha - 1} \exp\left[ -(\gamma - 1) x \right].
\label{eq:mhi_identity}
\end{align}
Using \eqref{eq:mhi_identity} with $\alpha = 0$,
$\int a_\gamma(x, y) \gamma \exp(-\gamma y) dy
= \gamma \exp( (1-\gamma ) x)
+ (1 - \gamma) \exp(-\gamma x)$.
So then
\eq{
&(\P_\gamma W^{\alpha})(x) - W^{\alpha}(x)
\\
&= 
\left\{ \frac{\gamma}{\alpha - 1} 
+ \frac{\gamma}{\gamma - \alpha}
+ \gamma - 1
\right\} 
\exp\left[ -(\gamma - \alpha) x \right] 
- \frac{\gamma}{\alpha - 1} \exp\left[ -(\gamma - 1) x \right]
- \gamma W^{\alpha}(x)^{\frac{1 + \alpha - \gamma}{\alpha}}.
}
It follows that
\begin{align}
(\P_\gamma W^{\alpha})(x) - W^{\alpha}(x)
&\le \frac{\gamma^*}{\alpha - 1} 
+ \frac{\gamma_*}{\gamma_* - \alpha}
+ \gamma^* - 1.
\label{eq:mhi_drift}
\end{align}
With the upper bound in \eqref{eq:mhi_drift}, we can now use Theorem~\ref{theorem:lb_with_target_tails} with $\phi \equiv c_*$, $\kappa = 1$, and $\alpha \in (1, \gamma_*)$.
We then have the lower bound for $\e$ and some $\delta_\e$:
\eq{
\inf_{\Q \in \S([0, \infty), (\gamma_*, \gamma^*)) } \inf_{\xi \in \C\left[ \A_{\Q}^{(t)}(\gamma_0, x_0, \cdot), \pi \right] } \xi( \{ x, y \in \X \times \X : |y - x| > \delta_{\e} \})
\ge \frac{
M_* (1 - \e)
}{ 
\left[ c_* t + \exp(\alpha x_0) \right]^{\frac{1}{\alpha - 1}} 
}
}
holds for every $t \in \Z_+$ uniformly in the adaptation strategy $\Q$.
\end{proof}

\end{document}
